\begin{scanpage}{0030}{13}
expansion converges only on part of the integration contour, the resultant series obtained after the integration will be semiconvergent (asymptotic). Limiting ourselves to the first term in it, we obtain:
\[\begin{aligned}v(r,\psi)&\approx\frac{\sin\frac\pi n}{\sqrt2\pi n}\frac{e^{i(kr+\frac\pi4)}}{\cos\frac\pi n-\cos\frac\psi n}\int_{-\infty}^{\infty}e^{-krs^2}\,ds=\\&=\frac{\frac1n\sin\frac\pi n}{\cos\frac\pi n-\cos\frac\psi n}\frac{e^{i(kr+\frac\pi4)}}{\sqrt{2\pi kr}}.\end{aligned}\tag{2.04}\]
The remaining terms of the asymptotic series have a value on the order of $\dfrac1{(kr)^{3/2}}$ and higher.

Expression (2.04) is valid with the condition $\left(\cos\frac\pi n-\cos\frac\psi n\right)\sqrt{kr}\gg1$ and describes that part of the diffraction field which has the character of cylindrical waves diverging from the wedge edge. With the incidence of the plane wave (1.23) on a wedge, the electric vector of which is parallel to the wedge edge, the cylindrical wave is determined in accordance with (1.25) and (2.04) by the equation
\[\begin{aligned}E_z=-H_\phi&=E_{0z}[v(r,\phi-\phi_0)-v(r,\phi+\phi_0)]=\\&=E_{0z}f\frac{e^{i(kr+\frac\pi4)}}{\sqrt{2\pi kr}},\end{aligned}\tag{2.05}\]
where
\[f=\frac{\sin\frac\pi n}{n}\left(\frac1{\cos\frac\pi n-\cos\frac{\phi-\phi_0}{n}}-\frac1{\cos\frac\pi n-\cos\frac{\phi+\phi_0}{n}}\right).\tag{2.06}\]
When the wedge is excited by the plane wave (1.24), in which the magnetic vector is parallel to the wedge edge, the cylindrical wave has the form
\[\begin{aligned}H_z=E_\phi&=H_{0z}[v(r,\phi-\phi_0)+v(r,\phi+\phi_0)]=\\&=H_{0z}g\frac{e^{i(kr+\frac\pi4)}}{\sqrt{2\pi kr}},\end{aligned}\tag{2.07}\]
\end{scanpage}
